\section{ANTECEDENTES}

El análisis territorial y temporal de la delincuencia ha evolucionado profundamente desde los primeros mapas de alfileres (\textit{pin maps}) utilizados por los departamentos de policía en el siglo XIX hasta las plataformas contemporáneas de inteligencia geoespacial asistidas por inteligencia artificial. Durante décadas, la criminología empírica descansó sobre postulados teóricos que explican la concentración delictiva en el espacio físico:

\begin{enumerate}
    \item \textbf{Teoría de las Actividades Rutinarias \cite{cohen1979}:} Propone que un delito ocurre cuando convergen en tiempo y espacio tres elementos fundamentales: un agresor motivado, un objetivo o víctima atractiva, y la ausencia de un guardián capaz. El movimiento cotidiano de la población define nodos de atracción delictiva.
    \item \textbf{Teoría del Patrón Delictivo \cite{brantingham1993}:} Establece que los infractores no eligen sus blancos de manera aleatoria, sino que exploran áreas familiares delimitadas por sus rutas de desplazamiento entre el hogar, el trabajo y las zonas de ocio, conformando corredores delictivos estructurados.
    \item \textbf{Criminología Ambiental y Ley de Concentración del Crimen \cite{weisburd2015}:} Evidencia empíricamente que un porcentaje diminuto de micro-unidades espaciales (segmentos de calle, esquinas o intersecciones) concentra más del 50\% de los eventos delictivos de una metrópoli completa.
\end{enumerate}

Con la digitalización de las administraciones públicas y el despliegue de iniciativas de datos abiertos (\textit{Open Data}), ciudades como Chicago han publicado catálogos históricos que superan los cientos de miles de registros del Departamento de Policía de Chicago (\textit{Chicago Police Department} -- CPD). Estos repositorios incluyen coordenadas geográficas de alta precisión, códigos unificados de incidentes (código IUCR de Illinois), sellos de tiempo, tipología de la locación y estatus de arresto. 

Sin embargo, en el contexto latinoamericano, y particularmente en el Estado Plurinacional de Bolivia, las bases de datos abiertas sobre criminalidad son escasas, fragmentadas o de acceso restringido. Ante esta deficiencia informativa, la ciudadanía y los analistas dependen primariamente del ecosistema de medios de comunicación y prensa digital (rotativos de alta circulación nacional y regional como \textit{El Deber}, \textit{Los Tiempos}, \textit{La Razón}, \textit{Correo del Sur}, \textit{El Potosí}, entre otros). La extracción de datos mediante raspado web (\textit{web scraping}) y el procesamiento de lenguaje natural (NLP) se transforman en una vía prioritaria para estructurar crónicas policiales y georreferenciar incidentes de seguridad ciudadana a escala nacional.

A nivel de técnicas computacionales, el estado del arte ha trascendido las limitaciones del análisis euclidiano bidimensional. En la superficie esférica de la Tierra, el uso indebido de distancias euclidianas sobre latitud y longitud proyectadas en grados introduce distorsiones geométricas críticas. El empleo sistemático de la distancia ortodrómica de \textit{Haversine}, sumado a la irrupción de algoritmos jerárquicos multiescala como HDBSCAN \cite{campello2013} y OPTICS \cite{ankerst1999}, ha posibilitado la detección simultánea de aglomeraciones densas en centros metropolitanos y concentraciones difusas en zonas periféricas. Complementariamente, los modelos neuronales profundos (\textit{Autoencoders}) han emergido como la técnica no supervisada de referencia para el aislamiento de anomalías espaciotemporales extremas.

